% Actual class: 03
% Slide decks: 3_neuron_layers; L3_summary
% Exact scope: 3_neuron_layers pp. 1--60; L3_summary pp. 1--31
% Video supplement: Lecture3_Part1 and Lecture3_Part2; durations confirmed, no standalone subtitle track
% Human verified: Yes

\section{第 03 节课：Neuron Layers}
\label{lecture:SC4001-L03}

\lecturemeta
  {SC4001-L03}
  {2026-08-28}
  {3\_neuron\_layers；L3\_summary}
  {主讲义 pp.~1--60；Summary pp.~1--31}
  {Lecture3\_Part1（40:37）与 Lecture3\_Part2（41:50）；视频时长已核对，公式与例题以课件为准}
  {已确认（2026-08-28；检查题完成，仅订正 $U_{5,2}$ 的术语）}

\subsection{本讲主线}

Neuron layer（神经元层）把 $K$ 个 neurons（神经元）的 affine computation（仿射计算）合并为 matrix operations（矩阵运算）。不同 layer 的 forward function（前向函数）和 local error signal（局部误差信号）不同，但参数 gradient 具有统一结构：
\[
  \boxed{\nabla_WJ=X^{\mathsf T}\nabla_UJ},
  \qquad
  \boxed{\nabla_{\bm b}J=(\nabla_UJ)^{\mathsf T}\bm 1_P}.
\]
这正是后续 multilayer backpropagation（多层反向传播）的基本模块。

\lecturetopic{SC4001-L03-T01}{Weight Matrix 与 Batch Forward Computation}

设 input dimension（输入维数）为 $n$，一层含 $K$ 个 neurons。将第 $k$ 个 neuron 的 weight vector（权重向量）$\bm w_k\in\mathbb R^n$ 按列排列：
\[
  W=[\bm w_1\ \bm w_2\ \cdots\ \bm w_K]\in\mathbb R^{n\times K},
  \qquad
  \bm b=(b_1,\ldots,b_K)^{\mathsf T}\in\mathbb R^K.
\]
对单个 input（输入）$\bm x\in\mathbb R^n$，
\[
  \boxed{\bm u=W^{\mathsf T}\bm x+\bm b}\in\mathbb R^K,
  \qquad
  \bm y=f(\bm u).
\]
若 $P$ 个 inputs 按行组成 $X\in\mathbb R^{P\times n}$，则
\[
  \boxed{U=XW+B=XW+\bm 1_P\bm b^{\mathsf T}}
  \in\mathbb R^{P\times K},
  \qquad Y=f(U).
\]
$U_{pk}$ 是第 $p$ 个 sample（样本）送入第 $k$ 个 neuron 后的 synaptic input / pre-activation（激活前输入）；真正的 output（输出）是 $Y_{pk}=f(U_{pk})$。因此本课程矩阵约定可记为：\textbf{row（行）对应 sample，column（列）对应 neuron}。

\begin{figure}[htbp]
  \centering
  \resizebox{0.98\textwidth}{!}{% Original diagram based on the matrix relationships in Lecture 3.
\begin{tikzpicture}[
  >=Latex,
  block/.style={draw=noteBlue,very thick,rounded corners=3pt,fill=blue!7,
    minimum width=2.55cm,minimum height=1.0cm,align=center},
  loss/.style={draw=ntuRed,very thick,rounded corners=3pt,fill=red!6,
    minimum width=1.7cm,minimum height=1.0cm,align=center},
  grad/.style={draw=verifyOrange,thick,rounded corners=3pt,fill=orange!8,
    minimum width=3.6cm,minimum height=0.9cm,align=center},
  flow/.style={->,thick,draw=noteBlue},
  back/.style={->,thick,draw=ntuRed}
]
  \node[block] (x) at (0,0) {$X\in\mathbb R^{P\times n}$\\batch inputs};
  \node[block] (u) at (3.7,0) {$U=XW+\bm 1_P\bm b^{\mathsf T}$\\$U\in\mathbb R^{P\times K}$};
  \node[block] (y) at (7.4,0) {$Y=f(U)$\\$Y\in\mathbb R^{P\times K}$};
  \node[loss] (j) at (10.35,0) {$J$\\loss};

  \draw[flow] (x) -- (u);
  \draw[flow] (u) -- (y);
  \draw[flow] (y) -- (j);

  \node[grad] (g) at (7.4,-2.0) {$G=\nabla_UJ$\\由 activation 与 loss 决定};
  \node[grad] (gw) at (2.4,-2.0) {$\nabla_WJ=X^{\mathsf T}G\in\mathbb R^{n\times K}$};
  \node[grad] (gb) at (2.4,-3.25) {$\nabla_{\bm b}J=G^{\mathsf T}\bm 1_P\in\mathbb R^K$};

  \draw[back] (j.south) |- node[pos=0.36,right,font=\scriptsize]{backpropagation} (g.east);
  \draw[back] (g.west) -- (gw.east);
  \draw[back] (g.south west) -- (gb.east);
  \draw[back,dashed] (x.south) |- (gw.north west);
\end{tikzpicture}
}
  \caption{Neuron layer 的 batch forward computation 与统一 gradient 结构。橙色框中的 $G=\nabla_UJ$ 由具体 activation（激活函数）和 loss（损失函数）决定。}
  \label{fig:sc4001-l03-layer-computation}
\end{figure}

\textbf{对应例题：}\relatedexercise{SC4001-L03-E01}{Lecture Example 1：Perceptron Layer Forward Pass}。

\lecturetopic{SC4001-L03-T02}{Single-layer Backpropagation、GD 与 SGD}

对 single sample（单样本），第 $k$ 个 neuron 满足
\[
  u_k=\bm w_k^{\mathsf T}\bm x+b_k,
  \qquad
  \nabla_{\bm w_k}J=\bm x\frac{\partial J}{\partial u_k}.
\]
将所有 neurons 合并，可得 outer product（外积）形式
\[
  \boxed{\nabla_WJ=\bm x(\nabla_{\bm u}J)^{\mathsf T}},
  \qquad
  \boxed{\nabla_{\bm b}J=\nabla_{\bm u}J}.
\]
所以课程定义的 SGD（每个 sample 更新一次）为
\[
  W\leftarrow W-\alpha\bm x(\nabla_{\bm u}J)^{\mathsf T},
  \qquad
  \bm b\leftarrow\bm b-\alpha\nabla_{\bm u}J.
\]
对 full-batch GD（全批量梯度下降），若 total loss（总损失）$J=\sum_{p=1}^{P}J_p$，则
\[
  \boxed{\nabla_WJ=X^{\mathsf T}\nabla_UJ},
  \qquad
  \boxed{\nabla_{\bm b}J=(\nabla_UJ)^{\mathsf T}\bm 1_P}.
\]
若实现使用 mean loss（平均损失），上述 batch gradients 还要除以 $P$；这会改变 gradient scale（梯度尺度），不能脱离 loss normalization（损失归一化）直接比较 learning rate（学习率）。

\begin{center}
\small
\begin{tabularx}{\textwidth}{@{}>{\raggedright\arraybackslash}p{3.0cm}>{\raggedright\arraybackslash}p{3.5cm}X@{}}
  \toprule
  Layer & Forward / loss & Local error signal $\nabla_UJ$ \\
  \midrule
  Linear layer（线性层） & $Y=U$；square error & $Y-D$ \\
  Perceptron layer（感知机层） & $Y=f(U)$；square error & $(Y-D)\odot f'(U)$ \\
  Softmax layer & $P=\operatorname{softmax}(U)$；cross-entropy & $P-T$ \\
  \bottomrule
\end{tabularx}
\end{center}

\lecturetopic{SC4001-L03-T03}{Perceptron Layer 与 Square-error Gradient}

按课程 terminology（术语），perceptron layer 使用连续 sigmoid activation（Sigmoid 激活）并完成 multidimensional nonlinear regression（多维非线性回归）：
\[
  y_k=f(u_k)=\frac{1}{1+e^{-u_k}},
  \qquad
  J=\frac12\sum_{k=1}^{K}(d_k-y_k)^2.
\]
Chain rule（链式法则）给出
\[
  \boxed{\nabla_{\bm u}J
    =-(\bm d-\bm y)\odot f'(\bm u)
    =(\bm y-\bm d)\odot f'(\bm u)}.
\]
对 batch，
\[
  \boxed{\nabla_UJ=(Y-D)\odot f'(U)}.
\]
Sigmoid 进入 saturation region（饱和区）时 $f'(u)$ 很小，即使 prediction error（预测误差）明显，传回的 gradient 也可能很弱。单层 sigmoid 虽是 nonlinear mapping（非线性映射），但每个 output 仍属于 $f(\bm w^{\mathsf T}\bm x+b)$ 这一受限函数族。

Summary 的 worked example（讲解例题）使用 scaled sigmoid
\[
  y=2\sigma(u)=\frac{2}{1+e^{-u}},
  \qquad
  \boxed{f'(u)=y\left(1-\frac y2\right)},
\]
使 output range（输出范围）变为 $(0,2)$。GD 每个 epoch 用全部 samples 更新一次；SGD 按顺序逐个处理 samples，后一个 sample 使用的是前一个 sample 已更新后的参数。

\textbf{对应例题：}\relatedexercise{SC4001-L03-E02}{Worked Example：Perceptron Layer 的 GD 与 SGD}。

\lecturetopic{SC4001-L03-T04}{Softmax、Cross-entropy 与 Logit Gradient}

Softmax layer 用于 $K$-class classification（$K$ 类分类）。对 logits（未归一化得分）$u_k$，
\[
  \boxed{p_k=P(y=k\mid\bm x)
    =\frac{e^{u_k}}{\sum_{j=1}^{K}e^{u_j}}},
  \qquad
  \hat y=\arg\max_k p_k=\arg\max_k u_k.
\]
各 $p_k$ 相互耦合且满足 $\sum_kp_k=1$。若 target（目标）为 one-hot vector（独热向量）$\bm t$，multiclass cross-entropy（多类交叉熵）为
\[
  J=-\sum_{k=1}^{K}t_k\ln p_k=-\ln p_d,
\]
其中 $d$ 是正确类别。利用 softmax Jacobian（雅可比矩阵）
\[
  \frac{\partial p_i}{\partial u_j}=p_i(\delta_{ij}-p_j),
\]
可得本讲最关键的简化：
\[
  \boxed{\nabla_{\bm u}J=\bm p-\bm t},
  \qquad
  \boxed{\nabla_UJ=P-T}.
\]
对 correct class（正确类别）$d$，$p_d-1<0$，gradient descent 会增大 $u_d$；对 incorrect class（错误类别），gradient 为 $p_k>0$，更新会减小相应 logit。

数值实现应先使用 shift invariance（平移不变性）
\[
  \operatorname{softmax}(\bm u)
  =\operatorname{softmax}(\bm u-c\bm 1),
  \qquad c=\max_k u_k,
\]
再计算指数，避免 numerical overflow（数值溢出）。

\lecturetopic{SC4001-L03-T05}{Softmax Decision Boundaries 与三分类例题}

类别 $i$ 与 $j$ 的 decision boundary（决策边界）满足
\begin{align*}
  p_i=p_j
  &\iff e^{u_i}=e^{u_j}
   \iff u_i=u_j\\
  &\iff
  \boxed{(\bm w_i-\bm w_j)^{\mathsf T}\bm x+(b_i-b_j)=0}.
\end{align*}
因此每一条 pairwise boundary（两两类别边界）都是 linear hyperplane（线性超平面）。多个类别的 regions（区域）由若干 hyperplanes 的交集形成，整体边缘可以是 piecewise-linear（分段线性）的，但单独一个 softmax layer 不能直接产生 curved boundary（弯曲边界）。Hidden layers（隐藏层）需要先学习 nonlinear features（非线性特征），再由 softmax 完成线性分类。

\begin{figure}[htbp]
  \centering
  \resizebox{0.70\textwidth}{!}{% Original reconstruction from the converged parameters in Lecture 3, Example 2.
\begin{tikzpicture}[>=Latex,scale=0.9]
  \draw[->,thick] (-5.1,0) -- (5.25,0) node[right] {$x_1$};
  \draw[->,thick] (0,-4.2) -- (0,4.35) node[above] {$x_2$};

  \begin{scope}
    \clip (-5,-4) rectangle (5,4);
    \draw[very thick,draw=ntuRed,domain=-5:5,samples=2]
      plot (\x,{2.167-0.517*\x});
    \draw[very thick,draw=noteBlue,domain=-5:5,samples=2]
      plot (\x,{2.425*\x-1.211});
    \draw[very thick,draw=verifyOrange,domain=-5:5,samples=2]
      plot (\x,{1.044+0.460*\x});
  \end{scope}

  \foreach \p in {(-1.0,3.0),(0.0,3.5),(2.0,3.0)}
    \draw[noteBlue,very thick,fill=blue!15] \p -- ++(-0.10,-0.18) -- ++(0.20,0) -- cycle;
  \foreach \p in {(-3.2,1.0),(-2.4,0.2),(-1.2,1.0)}
    \fill[ntuRed] \p circle (3pt);
  \foreach \p in {(2.0,0.1),(3.0,1.0),(4.0,0.1)}
    \draw[verifyOrange,very thick] \p ++(-0.10,-0.10) -- ++(0.20,0.20)
      \p ++(-0.10,0.10) -- ++(0.20,-0.20);

  \fill[black] (1.15,1.57) circle (2.2pt)
    node[above right,font=\scriptsize] {$u_A=u_B=u_C$};
  \node[noteBlue,font=\small] at (-0.5,3.65) {class A};
  \node[ntuRed,font=\small] at (-3.8,1.55) {class B};
  \node[verifyOrange,font=\small] at (3.8,1.55) {class C};

  \node[ntuRed,font=\scriptsize,rotate=-27] at (-2.9,3.3) {$u_A=u_B$};
  \node[noteBlue,font=\scriptsize,rotate=67] at (-0.7,-3.0) {$u_B=u_C$};
  \node[verifyOrange,font=\scriptsize,rotate=25] at (3.7,2.5) {$u_A=u_C$};
\end{tikzpicture}
}
  \caption{由 Lecture Example 2 收敛参数重绘的 three-class softmax boundaries。三条直线分别满足 $u_A=u_B$、$u_B=u_C$ 与 $u_A=u_C$；样本点仅作类别位置示意。}
  \label{fig:sc4001-l03-softmax-boundaries}
\end{figure}

讲义三分类例题的 initial forward pass（初始前向计算）得到 cross-entropy $34.36$、classification errors（分类错误数）$14$；一次 GD update 后参数明显改变。收敛时 classification error 为 $0$，但 cross-entropy 仍为 $0.58$：正确分类只要求正确类别 probability 最大，而 cross-entropy 还要求该 probability 尽量接近 $1$。

\textbf{对应例题：}\relatedexercise{SC4001-L03-E03}{Lecture Example 2：GD of a Softmax Layer}。

\FloatBarrier
\lecturetopic{SC4001-L03-T06}{Weight Initialization}

所有 neurons 使用完全相同的初始参数会产生相同 outputs 与 gradients，后续更新也保持相同，形成 symmetry problem（对称性问题）。但 weights 过大又会使 sigmoid/tanh 容易饱和。因此 initialization（初始化）应随机打破对称性，同时控制 forward activations 与 backward gradients 的 variance（方差）。

Xavier/Glorot uniform initialization（Xavier/Glorot 均匀初始化）为
\[
  w\sim U(-a,a),
  \qquad
  \boxed{a=\mathrm{gain}\sqrt{\frac{6}{n_{\mathrm{in}}+n_{\mathrm{out}}}}}.
\]
讲义列出的 gain（增益）为：linear 与 sigmoid 取 $1$，tanh 取 $5/3$，ReLU 取 $\sqrt2$，Leaky ReLU 取 $\sqrt{2/(1+\text{slope}^2)}$。

对 ReLU-family activations（ReLU 类激活），常用 Kaiming/He scaling（Kaiming/He 方差缩放）
\[
  \operatorname{std}(w)=\frac{\mathrm{gain}}{\sqrt{n_{\mathrm{in}}}}.
\]
\textbf{严格性补充：}课件把截断在 $2$ 个 standard deviations（标准差）内的 normal sampling（正态采样）直接称为 Kaiming normal；实际框架中的 Kaiming initialization 核心是按 fan-in/fan-out 和 nonlinearity 缩放方差，不一定采用 truncated normal（截断正态分布）。考试按课程公式作答，编程时应检查框架的具体定义。

\subsection{一分钟回顾}

Neuron layer 用 $U=XW+\bm1_P\bm b^{\mathsf T}$ 完成 batch forward computation；一旦得到 local error signal $\nabla_UJ$，所有 layer 都共享 $\nabla_WJ=X^{\mathsf T}\nabla_UJ$ 与 $\nabla_{\bm b}J=(\nabla_UJ)^{\mathsf T}\bm1_P$。Perceptron--square-error 给出 $(Y-D)\odot f'(U)$，softmax--cross-entropy 给出更简洁的 $P-T$。单独 softmax layer 的两两边界仍是 linear hyperplanes；合理 initialization 则负责打破对称性并维持信号尺度。
